\documentclass[12pt,a4paper]{article}
\usepackage[utf8]{inputenc}
\usepackage[portuguese]{babel}
\usepackage{hyperref}

\title{Eurovisao Napa 2025}
\author{Tomás Pinto - Scripting no Processamento de Linguagem Natural}
\date{\today}

\begin{document}

\maketitle

\begin{abstract}
Com base no modelo de n-grams desenvolvido e no método de scoring aplicado, as três frases selecionadas do texto são:
\begin{itemize}
    \item E o mais bonito tem sido o que acontece longe das câmaras, ou seja, conversas reais sobre música, países, histórias de vida.
    \item Nascidos numa cave na ilha da Madeira, os NAPA formaram-se entre amigos que partilhavam a paixão pela música.
    \item Numa conversa descontraída, entre gargalhadas cúmplices e reflexões honestas, os NAPA abriram o coração sobre a reação intensa que se seguiu à vitória no Festival da Canção — “foi duro, mas sabíamos ao que vínhamos”, admitem.
\end{itemize}
\end{abstract}

\section{Texto Original com Entidades (NER)}
Eurovisão 2025 | \textbf{Entrevista} (LOC) aos NAPA, representantes de \textbf{Portugal} (LOC) na \textbf{Suíça} (LOC)



A \textbf{Magazine} (ORG).HD esteve à conversa com os \textbf{NAPA} (ORG) , os representantes de \textbf{Portugal} (LOC) na \textbf{Eurovisão 2025} (MISC) , que este ano acontece em \textbf{Basileia} (LOC), na \textbf{Suíça} (LOC). A vitória no \textbf{Festival da Canção} (PER) foi inesperada e, para muitos, controversa -- a escolha da banda madeirense gerou divisões nas redes sociais, com elogios à sua autenticidade e críticas ao sistema de voto e o alegado afastamento dos padrões da música eurovisiva. Com o tempo, a poeira parece ter assentado, os ânimos acalmaram e o interesse em torno da banda deu lugar a uma expectativa palpável.



Para quem continua a dizer que esta não é uma canção "eurovisiva", vale a pena recordar que a \textbf{Eurovisão} (LOC) nunca foi sobre fórmulas fixas. O palco europeu é sinónimo de diversidade, onde cabem todos os géneros, todas as vozes e todas as emoções, independentemente dos resultados.



A proposta dos \textbf{NAPA} (ORG) é delicada, melancólica e honesta, como tantas outras que ali passaram em edições anteriores. Mais do que uma competição, a \textbf{Eurovisão} (LOC) é uma celebração da identidade de cada país e, nesse sentido, a canção que \textbf{Portugal} (LOC) leva em 2025 é profundamente nossa.



"Deslocado" é essencialmente uma história de migração e pertença. Nascidos numa cave na \textbf{ilha da Madeira} (LOC), os \textbf{NAPA} (ORG) formaram-se entre amigos que partilhavam a paixão pela música.



Carregam no som e nas letras as influências de \textbf{Arctic Monkeys} (ORG) , \textbf{Red Hot Chili Peppers} (ORG) , \textbf{The Beatles} (ORG) , \textbf{Caetano Veloso} (PER) e \textbf{Tom Jobim} (PER) -- uma mescla improvável que resulta numa identidade muito própria. O primeiro álbum, \textbf{" Senso Comum "} (MISC), chegou em 2019 e apresentou ao público a sua veia melódica e espírito introspectivo.



Em 2023, lançaram "Logo Se Vê", um disco mais maduro e seguro. Depois da primeira digressão nacional em 2024, a banda prepara agora o seu terceiro trabalho para o final de 2025. Mas, antes disso, o grande palco da \textbf{Basileia} (LOC) chama por eles.



\textbf{Guilherme Gomes} (PER) , \textbf{Lourenço Gomes} (PER) , \textbf{Francisco Sousa} (PER) , \textbf{Diogo Góis} (PER) e \textbf{João Rodrigues} (PER) falaram connosco sobre esta nova fase: as críticas que enfrentaram, o apoio que tem vindo a crescer e a canção com que querem tocar a \textbf{Europa} (LOC).



Numa conversa descontraída, entre gargalhadas cúmplices e reflexões honestas, os \textbf{NAPA} (ORG) abriram o coração sobre a reação intensa que se seguiu à vitória no \textbf{Festival da Canção} (PER) -- "foi duro, mas sabíamos ao que vínhamos\textbf{"} (MISC), admitem. Falaram da liberdade criativa que sempre defenderam desde os seus inícios, e de como transformaram a controvérsia numa força mobilizadora. Uma coisa é certa: os NAPA vão à \textbf{Eurovisão} (LOC) de cabeça erguida.



MHD: Será que nos poderiam revelar o processo criativo por detrás desta canção?



NAPA: O processo foi… nós convidámo-nos a nós próprios para fazer esta canção para o \textbf{Festival da Canção} (PER) . E eu já tinha esta música na gaveta há algum tempo. Ia para outro projeto, um álbum só de músicas sobre a \textbf{Madeira} (LOC). Mas depois percebi que seria ideal levá-la ao festival, porque conta a nossa história. Somos todos da \textbf{Madeira} (LOC), e aos 18 anos fomos estudar para \textbf{Lisboa} (LOC). Esse sentimento de adaptação a uma cidade nova, longe da família, foi sempre um desafio. Estive anos à procura das palavras certas para descrever isso, e acho que finalmente consegui.



A música ultrapassou a história pessoal e tornou-se um hino nacional. Como é que vêm esta transmutação da música, com a qual muitos agora se identificam?



A música conta uma história pessoal, mas é uma história que está sempre a acontecer. Apanhámos um pedacinho -- cinco rapazes da \textbf{Madeira} (LOC) -- mas há séculos que há portugueses a sair para procurar novas oportunidades. Hoje em dia continua a ser assim, principalmente por causa das dificuldades cá. E também sentimos que quanto mais pessoal é o que escreves, mais criativo se torna. Às vezes, se falares do teu jardinzinho, toda a gente quer saber como é esse jardim -- e ele pode ser mais parecido com o dos outros do que pensa.



Como têm corrido os ensaios? Há alguma coisa que estão a ajustar? E o que podemos esperar no dia 13 de maio?



Os ensaios aqui na \textbf{Basileia} (LOC) têm corrido bem. No primeiro dia foi aquele impacto de estar num palco gigante, com luzes e um \textbf{LED} (ORG) ball até ao chão. Mas sentimo-nos bem. A performance está como idealizámos, com pequenos ajustes. A verdade é que temos vindo a preparar isto há mais de dois meses. Agora, ver tudo materializado, rever os vídeos, ver como tudo encaixa no que imaginámos -- tem sido mesmo fixe.



E como tem sido o contacto com outros artistas da \textbf{Eurovisão?} (MISC) Já eram fãs ou é tudo novo para vocês?



Estávamos um bocado a leste deste mundo. \textbf{Claro} (ORG) que sabíamos do impacto da \textbf{Eurovisão} (LOC) e da vitória do \textbf{Salvador Sobral} (LOC), mas não conhecíamos a dimensão que isto tem. Entretanto, já conhecemos vários artistas nas pre-parties de \textbf{Madrid} (LOC) e \textbf{Amesterdão} (LOC). E o mais bonito tem sido o que acontece longe das câmaras, ou seja, conversas reais sobre música, países, histórias de vida. Essa ligação tem sido das coisas mais marcantes nesta experiência.



Sei que já têm concertos marcados para os próximos meses. Mas para além disso, há planos para um álbum ou colaborações com outros artistas portugueses?



\textbf{Calma} (PER)! \textbf{Calma!} (MISC) (risos) \textbf{Sim} (PER), já temos um álbum planeado. Ainda falta gravá-lo, mas as músicas estão feitas. E temos uma canção pronta para sair logo depois da \textbf{Eurovisão} (LOC). Não podemos revelar muito, mas vai ser bom. Queremos aproveitar esta exposição e lançar coisas rapidamente. Temos material guardado mesmo a pensar neste momento.



E quanto às vossas influências musicais? Que artistas vos inspiram?



Entre nós temos gostos muito semelhantes. A nível de banda , fomos influenciados por \textbf{Arctic Monkeys} (ORG), \textbf{Radiohead} (ORG), \textbf{Caetano Veloso} (PER), \textbf{Chico Buarque} (PER), \textbf{Los Hermanos} (ORG)… Também \textbf{Bon Iver} (ORG), \textbf{Mac DeMarco} (PER), \textbf{os Beatles} (ORG). E os \textbf{Bala Desejo} (ORG), uma banda brasileira mais recente. \textbf{Em Portugal} (LOC), claro que os \textbf{Capitão Fausto} (PER) são uma grande referência para escrevermos em português. Podemos dizer que também temos uma referência visual da artista madeirense \textbf{Lourdes de Castro} (PER). Foi ela que inspirou o visual da nossa performance. As roupas são da \textbf{Constança Entrudo} (PER), também ela madeirense.



Que mensagem gostariam de deixar aos portugueses que estão fora e que vão assistir à vossa performance?



A todos os portugueses lá fora -- se gostam da nossa música, se se identificam com a mensagem, votem no dia 13 de maio. Vamos representar todos os deslocados. \textbf{Percebemos} (LOC), ao partilhar esta música, que somos muitos. Muitos deslocados, portugueses e não só. E sentimos mesmo essa força coletiva.

\vspace{1cm}
\begin{thebibliography}{9}
\bibitem{eurovisao_napa_2025}
    Fonte analisada: \texttt{eurovisao\_napa\_2025}. Texto processado automaticamente.
\end{thebibliography}

\end{document}
